\section{Scope of the Problem}
Example~1 is one of a family of programs that Java accepts with no warning and that fail at run time.
This section first shows more programs of this family.

\subsection{More Failing Examples}
In Example~1, the local class \Q@MyFoo@ mentions \Q@X@ in its supertype \Q@Foo<X>@.
Example~2 exercises the same fault without needing a generic supertype.
Here \Q@MyFoo@ extends \Q@Object@, and \Q@X@ occurs only in the return type of the method \Q@x()@ of \Q@MyFoo@.
The method \Q@m@ returns \Q@Object@, so no type outside \Q@m@ mentions \Q@X@.

\noindent\begin{minipage}{\linewidth}
\begin{lstlisting}[numbers=left,xleftmargin=2em]
// Example 2
public class Broken {
  static <X> Object m(Object other, final X x, Function<X,String> f){
    class MyFoo{ public X x(){ return x; } }
    if(!(other instanceof MyFoo)){ return new MyFoo(); }
    MyFoo local= (MyFoo)other; // cast passes with no warnings
    System.out.println("old one was "+f.apply(local.x())); // throws ClassCastException
    return local;
  }
  public static void main(String[] a) {
    Object r1= Broken.<Integer>m("anything",5,i->""+i);
    Broken.<String>m(r1,"Hi",s->s);
  }
}
\end{lstlisting}
\end{minipage}

The first call of \Q@m@ instantiates \Q@X@ with \Q@Integer@, and returns a \Q@MyFoo@ whose method \Q@x()@ returns \Q@5@.
The second call of \Q@m@ instantiates \Q@X@ with \Q@String@.
In the second call, \Q@other instanceof MyFoo@ is true, and the cast \Q@(MyFoo)other@ succeeds with no warning.
Then \Q@local.x()@ has static type \Q@String@, but returns the \Q@Integer@ \Q@5@, and \Q@f.apply(local.x())@ throws a \Q@ClassCastException@.
The \Q@ClassCastException@ comes from the class that \Q@LambdaMetafactory@ generates at run time for the lambda \Q@s->s@.
javac compiles \Q@s->s@ into a synthetic method whose parameter has type \Q@String@,
and the generated \Q@apply(Object)@ casts its argument to \Q@String@ before calling that method.
The generated class is hidden,
%what is the generated class, make it more clear.
so its frame
%frames are for methods not classes.
does not appear in the default stack trace,
%is there a non default stack trace??
and the class file of \Q@Broken@ contains no \Q@checkcast@ for this cast.
%what do you mean? are you saying that a specific bytecode instruction is not present? if so we need to say it in this way and explain why it is relevant.

The program does not need to write the cast.
%the program does not write a cast, a programmer does.


%make a subsection for each example, or some other kind of divisor.
Example~3 exercises the same fault without needing an explicit cast.
Java accepts Example~3 with no warning.
Here we rely on the \Q@other instanceof MyFoo local@ test to bind
\Q@local@ in one step.

\noindent\begin{minipage}{\linewidth}
\begin{lstlisting}[numbers=left,xleftmargin=2em]
// Example 3
public class Broken {
  static <X> Object m(Object other, final X x, Function<X,String> f){
    class MyFoo{ public X x(){ return x; } }
    if(other instanceof MyFoo local){ // no cast, no warnings
      System.out.println("old one was "+f.apply(local.x())); // throws ClassCastException
    }
    return new MyFoo();
  }
  public static void main(String[] a) {
    Object r1= Broken.<Integer>m("anything",5,i->""+i);
    Broken.<String>m(r1,"Hi",s->s);
  }
}
\end{lstlisting}
\end{minipage}

Pattern matching in \Q@switch@ behaves in the same way.
In Example~4, the case \Q@case MyFoo local@ performs the same type test as the pattern of Example~3.

\noindent\begin{minipage}{\linewidth}
\begin{lstlisting}[numbers=left,xleftmargin=2em]
// Example 4
public class Broken {
  static <X> Object m(Object other, final X x, Function<X,String> f){
    class MyFoo{ public X x(){ return x; } }
    switch(other){
      case MyFoo local: // no cast, no warnings
        System.out.println("old one was "+f.apply(local.x())); // throws ClassCastException
        break;
      default:
        break;
    }
    return new MyFoo();
  }
  public static void main(String[] a) {
    Object r1= Broken.<Integer>m("anything",5,i->""+i);
    Broken.<String>m(r1,"Hi",s->s);
  }
}
\end{lstlisting}
\end{minipage}

The following listing is already covered by Examples~2 to~4.
The listing combines the \Q@switch@ pattern of Example~4 and the \Q@instanceof@ pattern of Example~3 in one method, and the \Q@switch@ fails first.
The listing uses \Q@Function<X,Integer>@ and the lambda \Q@s->s.length()@ instead of \Q@Function<X,String>@ and \Q@s->s@.
The failure is the failure of Examples~2 to~4: the class generated for the lambda casts its argument, the \Q@Integer@ \Q@5@, to \Q@String@ before calling the synthetic method \Q@lambda$main$1(String)@.
The class file of \Q@Broken@ has no \Q@checkcast@ to \Q@String@: the result of \Q@f.apply@ is used only in a string concatenation, and javac adds no \Q@checkcast@ to it.
In Examples~2 to~4, the class file of \Q@Broken@ does have a \Q@checkcast@ to \Q@String@, on the result of \Q@f.apply@, but that \Q@checkcast@ is never reached.
So in all four programs the failing cast is in the generated class, not in a \Q@checkcast@ of \Q@Broken@ and not in a virtual call.

\begin{lstlisting}
//Example that breaks in a NEW WAY in the bytecode
// https://godbolt.org/z/sMnxPG1TE
import java.util.function.Function;
class Broken {
  static <X> Object m(Object other, final X x, Function<X,Integer> f){
    class MyFoo{ public X x(){ return x; } }
    switch (other) {
        case MyFoo local:
            System.out.println("Pass with " + x); //snekly `local` is MyFoo<Integer> even when this is Brokem.<String>m(..)
            System.out.println("old one was "+f.apply(local.x()));//throws ClassCastException 
            break;
        default:
            break;
    }
    if(other instanceof MyFoo local){ //new instanceof+introduce var (pass with no warnings)
      System.out.println("Pass with " + x); //snekly `local` is MyFoo<Integer> even when this is Brokem.<String>m(..)
      System.out.println("old one was "+f.apply(local.x()));//throws ClassCastException    
     }
  return new MyFoo();//snekly it is MyFoo<Integer>
  }
  public static void main(String[] a) {
    Object r1= Broken.<Integer>m("anything",5,i->i);
    Broken.<String>m(r1,"Hi",s->s.length());
  }
}
// Above seems to fail but without explicit checkcast, failing one of the virtual method calls
\end{lstlisting}

\todo[inline]{The introduction promises forms that use class generics instead of method generics, and closures.}

\subsection{Generic Reification}
\todo[inline]{Not yet verified: the claims on Roslyn lifted classes, Rust (E0401, closure and opaque type identity, precise capturing), Swift, D, OCaml and SML local exceptions, and Scala outer tests.}
A natural suspicion is that this problem is ``simply'' a consequence of generic erasure:
a language that cannot distinguish the generic instantiations of a class at run time must have difficulties of this sort.
There is some truth to this suspicion: Java permits a cast that it cannot check, because of erasure.
However, Java does intend to flag uncheckable type tests:
for \Q@o@ of type \Q@Object@, Java flags the cast \Q@(List<String>)o@ with an unchecked warning, and rejects \Q@o instanceof List<String>@ with an error.
Java flags nothing in Examples~1 to~4 because \Q@MyFoo@ is not generic:
the type test of \Q@other@ against \Q@MyFoo@ has no generic argument to leave unchecked.

Reification does not change this by itself.
A reified runtime can only check what the run-time identity of a class contains.
If the run-time identity of \Q@MyFoo@ is just the declared class \Q@MyFoo@, reification adds nothing to the cast \Q@(MyFoo)other@.
Java treats the class \Q@MyFoo extends Foo<X>@ as it treats the class \Q@String extends Comparable<String>@: a cast inspects only the declared class.
A cast to \Q@String@ does not inspect the generic argument of \Q@Comparable<String>@,
and a cast to \Q@MyFoo@ does not inspect the generic argument of \Q@Foo<X>@.
But the argument \Q@String@ of \Q@Comparable<String>@ is a constant of the class \Q@String@: a reified runtime records it once, in the class table of \Q@String@.
The argument \Q@X@ of \Q@Foo<X>@ is not a constant of the class \Q@MyFoo@: \Q@X@ is free in the declaration of \Q@MyFoo@, and \Q@X@ changes between calls of \Q@m@.

The implementations of method-local classes, and of the similar features of other languages, differ in the granularity of the run-time identity of the local class.
There are three granularities: one identity per declared class, one identity per instantiation of the enclosing scope, and one identity per evaluation of the declaration.
Some languages avoid the choice by forbidding the capture of generic types.

\paragraph*{One identity per declared class.}
Java has this granularity, and so would a version of Java with reified generics that keeps \Q@MyFoo@ non-generic.
In Example~1, such a reified Java needs the argument of \Q@Foo<X>@ in each instance of \Q@MyFoo@, since a reified runtime knows the generic arguments of every instance of \Q@Foo@.
All instances of \Q@MyFoo@ share one class table, so each instance keeps a field with the argument of \Q@Foo@.
This leaves two possible implementations of the cast \Q@(MyFoo)other@:
\begin{enumerate}
\item The cast checks only the class table, and the cast succeeds. This is the unsoundness of Java, on a reified runtime.
\item The cast checks the class table and also the stored argument of \Q@Foo@, as a cast to \Q@Foo<Integer>@ would. This rejects Example~1, but only because \Q@X@ occurs in the supertype \Q@Foo<X>@.
\end{enumerate}
In Examples~2 to~4, \Q@MyFoo@ extends \Q@Object@: there is no stored argument to check, and the second implementation of the cast succeeds too.
The second implementation checks exactly the captured generic types that occur in the supertypes of the local class.
That is, the second implementation funnels only those captured generic types, and leaves the other captured generic types unchecked.
Moreover, in Examples~2 to~4 the signature of \Q@x()@ mentions \Q@X@, but \Q@MyFoo@ has no generic parameter to bind \Q@X@ in a reified runtime.
A reified Java must either erase \Q@X@ inside \Q@MyFoo@, or make \Q@MyFoo@ generic: that is, funnel \Q@X@.

\paragraph*{One identity per instantiation of the enclosing scope.}
In a language which implements parametric polymorphism by templating / heterogeneous compilation
(C++, or early experiments with generic Java~\cite{agesen1997adding}),
each instantiation of \Q@m@ with different generic parameters is compiled separately, resulting in separate
class declarations for each instantiation.
Consider reinterpreting Example~1 in C++.

\noindent\begin{minipage}{\linewidth}
\begin{lstlisting}[numbers=left,xleftmargin=2em,language=c++]
// Example 1 in C++
class Object{ public: virtual ~Object(){} };
template<typename X> class Foo : public Object {
    public: virtual X x() { return x(); } };
template<typename X>
Foo<X>* m(Object* other, const X* x) {
    class MyFoo : public Foo<X> {
        const X* v;
      public: MyFoo(const X* v):v(v){} X x() { return *v; } };
    MyFoo* d = dynamic_cast<MyFoo*>(other);(*\label{line:cpp_cast}*)
    if(d == nullptr) { return new MyFoo(x); }
    return d;
}
\end{lstlisting}
\end{minipage}

A C++ local class cannot use the parameters of its enclosing function, so \Q@MyFoo@ receives \Q@x@ through its constructor.
\Q@Object@ has a virtual destructor, since \Q@dynamic_cast@ requires a polymorphic source type.
The calls \Q@r1=m<int>(nullptr,&i)@ and then \Q@m<std::string>(r1,&s)@ mirror Example~1.
In the second call, the \Q@dynamic_cast@ at Line \ref{line:cpp_cast} returns \Q@nullptr@:
\Q@m<int>@ and \Q@m<std::string>@ declare two distinct classes \Q@MyFoo@, and \Q@r1@ is an instance of the first one.
Thus the second call returns a new \Q@MyFoo@, as a sound type test should.
In its diagnostics, \texttt{gcc} names the first class \Q@m<int>(Object*, const int*)::MyFoo@:
the instantiation of the method signature is the scope of the declaration of \Q@MyFoo@.
Inside the template \Q@m<X>@, the C++ type system sees \Q@MyFoo@ as the dependent type \Q@m<X>::MyFoo@.
So C++ does expose \Q@X@ in the type of \Q@MyFoo@: through the name of the enclosing scope instead of through a generic parameter of \Q@MyFoo@.
D works in the same way, and its ``Voldemort types'' make the scope visible:
a local struct \Q@S@ returned from the template function \Q@f@ has the type \Q@f!int.S@ when \Q@f@ is instantiated with \Q@int@~\tocite{D Voldemort types}.

C\# has no local classes, but its compiler lifts closures, local functions, iterators and async methods into generated classes.
Inside a generic method \Q@M<T>@, the generated classes are themselves generic in the captured \Q@T@, as in \Q@<>c__DisplayClass0_0<T>@~\tocite{Roslyn lowering}.
Since .NET reifies generics, the generated classes for \Q@M<int>@ and \Q@M<string>@ have distinct run-time identities.
These generated classes can not be named in C\# source, so no cast to them can be written.

Rust closures and \Q@impl Trait@ opaque types capture the generic parameters of the enclosing function,
and the type of the closure, or the opaque type, includes those generic parameters in its identity:
the \Q@TypeId@ of a closure created in \Q@m::<i32>@ differs from the \Q@TypeId@ of the same closure created in \Q@m::<String>@~\cite{rustClosures}.
The precise capturing syntax of Rust 2024, \Q@impl Trait + use<T>@, lists explicitly the generic parameters that an opaque type captures~\tocite{Rust precise capturing RFC 3617}.

Heterogeneous compilation has subtle costs and wide-ranging consequences for language design, not just implementation,
and complex performance consequences~\cite{odersky2001two,kennedy2001design}.
However, these costs are not intrinsic to this granularity:
.NET keeps a single implementation for the instantiations at reference types and still gives each instantiation its own run-time identity~\cite{kennedy2001design}.

Each of these sound designs makes the captured generic types part of the identity of the local class.
That is, each design exposes the captured generic types:
to the static type system and to the runtime (C++, D, Rust), or only to the runtime (C\# generated classes).
Each design is a form of funnelling, done by the compiler or by the runtime instead of in the source.

\paragraph*{One identity per evaluation of the declaration.}
In OCaml, \Q@let exception E of a in ...@ creates a fresh exception constructor each time the declaration is evaluated;
local exception declarations in Standard ML are generative in the same way~\tocite{OCaml local exceptions, SML generative exceptions}.
A handler for \Q@E@ does not catch an \Q@E@ raised by another evaluation of the declaration, whatever the instantiation of \Q@a@.
Scala has the same granularity for captured values instead of captured types.
An instance of an inner class \Q@Inner@ captures its enclosing instance \Q@$outer@, and its type is the path-dependent type \Q@a.Inner@.
For the pattern \Q@case i: a.Inner@, the Scala compiler emits an ``outer test'' that compares the captured \Q@$outer@ with \Q@a@~\tocite{Scala outer tests}.
Even here the type system exposes what is captured: inside the class of \Q@a@, \Q@Inner@ means \Q@this.Inner@.
This granularity is sound, but it is stricter than needed:
a \Q@MyFoo@ created by an earlier call of \Q@m@ with the same instantiation of \Q@X@ would also be rejected.
This changes the behaviour of programs that are correct with one identity per instantiation of the enclosing scope.

\paragraph*{Forbidden capture.}
In Rust, a struct, function or impl declared inside a generic function can not use the generic parameters of the enclosing function (error E0401).
The programmer must declare those generic parameters again on the inner item, which is funnelling done by hand.
Swift rejects a type declared inside a generic function~\tocite{Swift: type cannot be nested in generic function}.

\paragraph*{Capturing is exposing.}
We know of no language that keeps a captured generic type private to a local class (absent from the source, from the static type system and from the run-time identity of the class) and still checks the captured generic type at run time.
As soon as an implementation checks the captured generic type, the captured generic type is part of the identity of the local class: the local class is funnelled.

Funnelling is the static form of one identity per instantiation of the enclosing scope:
in Example~1, \Q@MyFoo@ funnels \Q@X@ and becomes \Q@MyFoo<X>@.
With erasure, the existing rule of Java then rejects \Q@other instanceof MyFoo<X>@ for \Q@other@ of type \Q@Object@, because \Q@Object@ does not determine \Q@X@.
These are the two compile-time points of the introduction.
With reification, the same type test returns false in the second call of \Q@m@.
These are the two run-time points of the introduction.
The same exposure of \Q@X@ enables all four points.
Thus funnelling is independent of how generics are implemented, and Section~\ref{sec:core} develops funnelling in a core language without reification.

\todo[inline]{ah, we need to check Flix, where maybe we can cook up something roughly like this?}

This is a problem in other languages (cite cite) as shown in greater details in Rel Work.

Now, assume we have a readable small step reduction over a version of FJ rich enough to express the problem.
We will assume: 
- Simple generics
- A syntax like anonymous inner classes, but where we give them names. We can call them
named inner classes.
This imports the concept of method local inner classes within an expression based calculus.
//our model will also need casts + instanceof
\begin{lstlisting}
class A{ A onlyA()-> new A() }
class B{ B onlyB()-> new B() }
class Foo<X> extends Object{ X x() {return x();} }
public class Broken {
  <X> Foo<X> m(Object other, X x) ->
    if(other instanceof MyFoo)
      then ((MyFoo)other)
      else new MyFoo() extends Foo<X>{ X x(){ return x; } }//Bad
  }
main:
  new Broken().<B>m(new Broken().<A>m(new Object(),new A()),new B()).x().onlyB();
\end{lstlisting}


\begin{lstlisting}
class Foo<X> extends Object{ X x() {return x();} }
public class Broken {
  <X> Foo<X> m(Object other, X x) ->
    if(other instanceof MyFoo<X>) //unchecked instanceof (ill typed in our core)?
      then ((MyFoo<X>)other)//unchecked cast (ill typed in our core)?
      else new MyFoo<X>() extends Foo<X>{ public X x(){ return x; } }//good
  }
\end{lstlisting}
\mjnote{answered by Figure~\ref{fig:cast}: with \texttt{other : Object}, \texttt{other instanceof MyFoo<X>} is ill typed, because \texttt{X} is not determined by \texttt{Object} (as javac: \texttt{Object cannot be safely cast to MyFoo<X>}). With \texttt{other : Foo<X>} it is well typed and sound under erasure.}

Note how funnelling purpose is to insert the generics. Without funnelling, MyFoo is a non generic type, so the question about generic reification does not even arise.

small step reduction
\begin{verbatim}
01:  new Broken().<B>m(new Broken().<A>m(new Object(),new A()),new B()).x().onlyB();
02:  new Broken().<B>m(
    if(new Object() instanceof MyFoo<X>)
      then (MyFoo<A>) new Object()
      else new MyFoo<A>() extends Foo<A>{ public A x(){ return new A(); } }
      ,new B()).x().onlyB();
03:  new Broken().<B>m(
    if(false)
      then (MyFoo<A>) new Object()
      else new MyFoo<A>() extends Foo<A>{ public A x(){ return new A(); } }
      ,new B()).x().onlyB();
04:  new Broken().<B>m(
       new MyFoo<A>() extends Foo<A>{ public A x(){ return new A(); } }
      ,new B()).x().onlyB();
05:   (if(new MyFoo<A>() extends Foo<A>{ public A x(){ return new A(); } } instanceof MyFoo<B>)
       then (MyFoo<A>) new MyFoo<A>() extends Foo<A>{ public A x()->new A() }
       else new MyFoo<B>() extends Foo<B>{ public B x(){ return new B(); } }).x().onlyB()
06a:  (if(false)//reified gen
       then (MyFoo<B>) new MyFoo<A>() extends Foo<A>{ public A x(){ return new A(); } }
       else new MyFoo<B>() extends Foo<B>{ public B x(){ return new B(); } }).x().onlyB()
07a:  (new MyFoo<B>() extends Foo<B>{ public B x(){ return new B(); } }).x().onlyB()
08a:  new B().onlyB()
09a:  new B()//success and different semantic
06b:  (if(true)//Java
       then (MyFoo<B>) new MyFoo<A>() extends Foo<A>{ public A x(){ return new A(); } }
       else new MyFoo<B>() extends Foo<B>{ public B x(){ return new B(); } }).x().onlyB()
07b:  ((MyFoo<B>) new MyFoo<A>() extends Foo<A>{ public A x(){ return new A(); } }).x().onlyB()
08b:  (new MyFoo<A>() extends Foo<A>{ public A x()->new A() }).x().onlyB()
     //assuming cast and instanceof behave consistently (unchecked)
09b:  new A().x().onlyB()
10b:  new A().onlyB() //stuck!!
6c:  (if(err)//compile or at least run time
       then ((MyFoo<B>) new MyFoo<A>() extends Foo<A>{ public A x(){ return new A(); } }
       else new MyFoo<B>() extends Foo<B>{ public B x(){ return new B(); } }).x().onlyB()
\end{verbatim}

// Maybe above we want MyFoo:Object<A> and MyFoo:Object<B>?

% interface Shape{ Square toSquare(); }
% class Square extends Shape{Square toSquare(){ return this; } }
%
%
%
% public class Broken {
%   static <X> Foo<X,?> m(Top<?> other, final X x){
%     class MyFoo extends Foo<X,MyFoo>{
%       public X x(){ return x; }
%       public MyFoo myFooOrElse(MyFoo myFoo){ return this; }
%     }
%     return other.selfOrElse(new MyFoo());
%   }
%   public static void main(String[] a) {
%     Foo<Integer> r1= Broken.<Integer>m(new Top<String>(),5);
%     Foo<String> r2= Broken.<String>m(r1,"Hi");
%     System.out.println(r1.x());//5
%     System.out.println(r2.x());//Throws ClassCastException: ..
%   }
% }
% class Top<S>{
%   public S selfOrElse(S s){ return s; }
%   }
% class Foo<X,S> extends Top<S>{ public X x() {return x();} }
%
%
